\chapter{DOGFOOD}\label{record-042}

\noindent{\footnotesize\color{BookMuted}\path{docs/strategy/dogfood.md}\par}\smallskip

\begin{quote}
Historical note: this milestone framing is kept for provenance. The current docs-surface signpost is \href{https://github.com/flyingrobots/bijou/blob/cbd70612140e7a450d764b3c7be1c6af878b99bc/docs/DOGFOOD.md}{\texttt{d\allowbreak{}o\allowbreak{}c\allowbreak{}s\allowbreak{}/\allowbreak{}D\allowbreak{}O\allowbreak{}G\allowbreak{}F\allowbreak{}O\allowbreak{}O\allowbreak{}D\allowbreak{}.\allowbreak{}m\allowbreak{}d}}, and new delivery work should start in METHOD backlog/design surfaces instead of here.
\end{quote}

\emph{Documentation Of Good Foundational Onboarding and Discovery}

\emph{Milestone strategy for the post-v4 living docs platform}

\section{Why this milestone exists}\label{record-042-why-this-milestone-exists}

After \texttt{v\allowbreak{}4\allowbreak{}.\allowbreak{}0\allowbreak{}.\allowbreak{}0}, the next problem is not ``what capability can Bijou add next?'' It is:

\begin{itemize}
\tightlist
\item
  how does someone understand Bijou quickly?
\item
  how do docs, demos, examples, and tests stop drifting apart?
\item
  how do we show graceful lowering instead of merely describing it?
\item
  how do we make the system easy to explore without multiplying one-off example apps?
\end{itemize}

This milestone is the answer. It is the first post-v4 product milestone because it improves understanding, trust, and development speed at the same time.

\texttt{D\allowbreak{}O\allowbreak{}G\allowbreak{}F\allowbreak{}O\allowbreak{}O\allowbreak{}D} is a better name than ``Milestone 1'' because it reflects the actual philosophy:

\begin{itemize}
\tightlist
\item
  docs should be experiential, not only descriptive
\item
  components should be explored in the medium they are built for
\item
  users should be able to compare states, modes, and patterns directly
\item
  the same source of truth should drive teaching, testing, and tooling
\end{itemize}

\section{Design-thinking frame}\label{record-042-design-thinking-frame}

\subsection{Primary user}\label{record-042-primary-user}

The primary user is the \textbf{curious builder}:

\begin{itemize}
\tightlist
\item
  a developer evaluating Bijou for a real terminal application
\item
  a contributor trying to understand a component family or pattern
\item
  a returning maintainer trying to remember what is canonical and what is not
\end{itemize}

They are not asking for infinite flexibility first. They are asking:

\begin{itemize}
\tightlist
\item
  ``What is this thing?''
\item
  ``When do I use it?''
\item
  ``What does it look like in real terminal conditions?''
\item
  ``How does it behave in \texttt{r\allowbreak{}i\allowbreak{}c\allowbreak{}h}, \texttt{s\allowbreak{}t\allowbreak{}a\allowbreak{}t\allowbreak{}i\allowbreak{}c}, \texttt{p\allowbreak{}i\allowbreak{}p\allowbreak{}e}, and \texttt{a\allowbreak{}c\allowbreak{}c\allowbreak{}e\allowbreak{}s\allowbreak{}s\allowbreak{}i\allowbreak{}b\allowbreak{}l\allowbreak{}e}?''
\item
  ``Can I trust that the example, docs, and tests all describe the same thing?''
\end{itemize}

\subsection{Secondary user}\label{record-042-secondary-user}

The secondary user is the \textbf{component author / framework maintainer}:

\begin{itemize}
\tightlist
\item
  they need a safe place to evolve a component family
\item
  they need stories and scenarios that can double as docs and regressions
\item
  they need a way to inspect behavior without hand-rolling a bespoke demo for every new idea
\end{itemize}

\subsection{Core jobs to be done}\label{record-042-core-jobs-to-be-done}

\begin{enumerate}
\def\labelenumi{\arabic{enumi}.}
\item
  \textbf{Help me understand Bijou without hunting across the repo.} I want one place where the doctrine, the live behavior, and the example source agree.
\item
  \textbf{Help me compare interaction profiles and states quickly.} I want to see the same component in \texttt{r\allowbreak{}i\allowbreak{}c\allowbreak{}h}, \texttt{s\allowbreak{}t\allowbreak{}a\allowbreak{}t\allowbreak{}i\allowbreak{}c}, \texttt{p\allowbreak{}i\allowbreak{}p\allowbreak{}e}, and \texttt{a\allowbreak{}c\allowbreak{}c\allowbreak{}e\allowbreak{}s\allowbreak{}s\allowbreak{}i\allowbreak{}b\allowbreak{}l\allowbreak{}e} without writing four custom demos.
\item
  \textbf{Help me trust what I am seeing.} I want docs, demos, and regression scenarios to come from one structured source of truth instead of diverging copies.
\item
  \textbf{Help me evolve the system safely.} I want to add or refine a component family once and have that work propagate into docs, previews, and test surfaces.
\end{enumerate}

\section{The current problem}\label{record-042-the-current-problem}

Today, Bijou is much healthier than it was pre-v4, but the learning surface is still fragmented:

\begin{itemize}
\tightlist
\item
  doctrine lives in \texttt{d\allowbreak{}o\allowbreak{}c\allowbreak{}s\allowbreak{}/\allowbreak{}d\allowbreak{}e\allowbreak{}s\allowbreak{}i\allowbreak{}g\allowbreak{}n\allowbreak{}-\allowbreak{}s\allowbreak{}y\allowbreak{}s\allowbreak{}t\allowbreak{}e\allowbreak{}m\allowbreak{}/\allowbreak{}*}
\item
  public teaching lives in package READMEs and example READMEs
\item
  interactive demonstrations live in example programs
\item
  regression coverage lives in smoke harnesses, PTY scripts, and targeted tests
\end{itemize}

Each of those serves a purpose, but they do not share one canonical representation of ``a component family in context.'' That fragmentation increases maintenance cost and makes the product harder to learn than it needs to be.

\section{Principles}\label{record-042-principles}

This milestone should follow these principles:

\begin{enumerate}
\def\labelenumi{\arabic{enumi}.}
\item
  \textbf{One source of truth} Stories should be structured records that can feed docs, demos, previews, and tests.
\item
  \textbf{Show behavior, not only prose} Lowering, width sensitivity, and interaction changes should be visible side-by-side.
\item
  \textbf{Docs are product} The docs experience is not a sidecar artifact. It is part of Bijou's value proposition.
\item
  \textbf{Terminal-native first} Bijou should document itself in its own medium where possible.
\item
  \textbf{Regression surfaces should emerge from teaching surfaces} If a story is canonical enough to teach, it should be close to canonical enough to test.
\end{enumerate}

\section{Scope of DOGFOOD}\label{record-042-scope-of-dogfood}

\subsection{Current validation slice}\label{record-042-current-validation-slice}

The first implementation slice now lives in \href{https://github.com/flyingrobots/bijou/blob/cbd70612140e7a450d764b3c7be1c6af878b99bc/examples/docs/README.md}{examples/docs}.

It is intentionally small:

\begin{itemize}
\tightlist
\item
  three canonical component stories
\item
  one framed docs shell
\item
  one docs pane derived from structured story fields
\item
  one live demo pane driven by the same story record
\item
  profile switching across \texttt{r\allowbreak{}i\allowbreak{}c\allowbreak{}h}, \texttt{s\allowbreak{}t\allowbreak{}a\allowbreak{}t\allowbreak{}i\allowbreak{}c}, \texttt{p\allowbreak{}i\allowbreak{}p\allowbreak{}e}, and \texttt{a\allowbreak{}c\allowbreak{}c\allowbreak{}e\allowbreak{}s\allowbreak{}s\allowbreak{}i\allowbreak{}b\allowbreak{}l\allowbreak{}e}
\end{itemize}

That slice is not the final docs product. It is the first proof that the story substrate can drive teaching and live previews without re-encoding the doctrine by hand.

It also exposed a broader shell need: when a docs lane is visually focused, that lane should own its local controls. The resulting shell/input design is captured in \href{https://github.com/flyingrobots/bijou/blob/cbd70612140e7a450d764b3c7be1c6af878b99bc/docs/strategy/focus-owns-input.md}{Focus Owns Input} and \href{https://github.com/flyingrobots/bijou/blob/cbd70612140e7a450d764b3c7be1c6af878b99bc/docs/specs/focused-pane-input.spec.json}{Focused Pane Input}.

The docs shell also now has two obvious shell-level proving surfaces:

\begin{itemize}
\tightlist
\item
  component search should feel like a standard shell affordance, not a bespoke docs shortcut
\item
  settings should open as a shell-owned drawer rather than as a custom page overlay
\end{itemize}

The resulting settings direction is captured in \href{https://github.com/flyingrobots/bijou/blob/cbd70612140e7a450d764b3c7be1c6af878b99bc/docs/strategy/settings-belong-to-the-shell.md}{Settings Belong to the Shell} and \href{https://github.com/flyingrobots/bijou/blob/cbd70612140e7a450d764b3c7be1c6af878b99bc/docs/specs/frame-settings-drawer.spec.json}{Frame Settings Drawer}.

As the settings drawer matured, another design truth became clear: if DOGFOOD is serious about proving Bijou, polished preference-row presentation should graduate into canonical Bijou componentry instead of remaining trapped inside shell-local drawing code. That direction is now captured in \href{https://github.com/flyingrobots/bijou/blob/cbd70612140e7a450d764b3c7be1c6af878b99bc/docs/strategy/preference-lists-belong-to-bijou.md}{Preference Lists Belong to Bijou} and \href{https://github.com/flyingrobots/bijou/blob/cbd70612140e7a450d764b3c7be1c6af878b99bc/docs/specs/preference-list-surface.spec.json}{Preference List Surface}.

As the docs app grew, that broadened into a larger shell-quality question: the title screen, header, footer, search, help, settings, notifications, confirmation flows, and quit semantics all need to feel like one intentional product. That direction is now captured in \href{https://github.com/flyingrobots/bijou/blob/cbd70612140e7a450d764b3c7be1c6af878b99bc/docs/strategy/humane-shell.md}{The Humane Shell}.

The next shell-feedback layer under that umbrella is notification review: transient notices need a standard home once they disappear. That direction is now captured in \href{https://github.com/flyingrobots/bijou/blob/cbd70612140e7a450d764b3c7be1c6af878b99bc/docs/strategy/notification-history-belongs-to-the-shell.md}{Notification History Belongs to the Shell} and \href{https://github.com/flyingrobots/bijou/blob/cbd70612140e7a450d764b3c7be1c6af878b99bc/docs/specs/shell-notification-center.spec.json}{Shell Notification Center}.

\subsection{1. Story Protocol}\label{record-042-story-protocol}

This is the highest-leverage foundation.

Reference documents:

\begin{itemize}
\tightlist
\item
  \href{https://github.com/flyingrobots/bijou/blob/cbd70612140e7a450d764b3c7be1c6af878b99bc/docs/strategy/component-story-v0.md}{ComponentStory v0 design note}
\item
  \href{https://github.com/flyingrobots/bijou/blob/cbd70612140e7a450d764b3c7be1c6af878b99bc/docs/specs/story-protocol.spec.json}{Story Protocol spec}
\end{itemize}

What it is:

\begin{itemize}
\tightlist
\item
  a structured \texttt{C\allowbreak{}o\allowbreak{}m\allowbreak{}p\allowbreak{}o\allowbreak{}n\allowbreak{}e\allowbreak{}n\allowbreak{}t\allowbreak{}S\allowbreak{}t\allowbreak{}o\allowbreak{}r\allowbreak{}y} / \texttt{P\allowbreak{}a\allowbreak{}t\allowbreak{}t\allowbreak{}e\allowbreak{}r\allowbreak{}n\allowbreak{}S\allowbreak{}t\allowbreak{}o\allowbreak{}r\allowbreak{}y} model
\item
  enough metadata to drive:

  \begin{itemize}
  \tightlist
  \item
    docs pages
  \item
    live demos
  \item
    examples/showcase entries
  \item
    scenario tests
  \item
    replay artifacts
  \end{itemize}
\end{itemize}

What it should capture:

\begin{itemize}
\tightlist
\item
  title and summary
\item
  family / pattern identity
\item
  use / avoid guidance
\item
  graceful lowering descriptions
\item
  demo inputs / states
\item
  width and profile presets
\item
  optional source snippet references
\item
  optional selectors / semantics
\item
  optional scenario hooks
\end{itemize}

Why it matters:

\begin{itemize}
\tightlist
\item
  without this, every later investment duplicates logic
\item
  with it, the docs TUI, scenario harness, lab mode, and replay tooling all share the same substrate
\end{itemize}

Primary assumption to test:

\begin{itemize}
\tightlist
\item
  one story format can serve both human-facing docs and machine-facing regression/tooling needs without becoming too heavy
\end{itemize}

\subsection{2. Bijou Docs TUI}\label{record-042-bijou-docs-tui}

This is the flagship user-facing product surface.

What it is:

\begin{itemize}
\tightlist
\item
  a Bijou-native docs application
\item
  left navigation, content pane, demo pane
\item
  width and interaction-profile controls
\item
  source and guidance next to live render output
\end{itemize}

Why it matters:

\begin{itemize}
\tightlist
\item
  it lets Bijou document itself in the medium it is designed for
\item
  it makes the design system feel real instead of abstract
\item
  it provides a natural home for component families, patterns, and canonical demos
\end{itemize}

Primary assumption to test:

\begin{itemize}
\tightlist
\item
  developers learn Bijou materially faster from a terminal-native docs surface than from disconnected README prose
\end{itemize}

\subsection{3. Interaction Profile Simulator}\label{record-042-interaction-profile-simulator}

This makes graceful lowering concrete.

What it is:

\begin{itemize}
\tightlist
\item
  a shared control surface for \texttt{r\allowbreak{}i\allowbreak{}c\allowbreak{}h}, \texttt{s\allowbreak{}t\allowbreak{}a\allowbreak{}t\allowbreak{}i\allowbreak{}c}, \texttt{p\allowbreak{}i\allowbreak{}p\allowbreak{}e}, \texttt{a\allowbreak{}c\allowbreak{}c\allowbreak{}e\allowbreak{}s\allowbreak{}s\allowbreak{}i\allowbreak{}b\allowbreak{}l\allowbreak{}e}
\item
  likely paired with width presets, reduced motion, and no-color modes
\end{itemize}

Why it matters:

\begin{itemize}
\tightlist
\item
  graceful lowering is one of Bijou's clearest differentiators
\item
  right now it is described better than it is demonstrated
\end{itemize}

Primary assumption to test:

\begin{itemize}
\tightlist
\item
  showing profile shifts directly will reduce both documentation ambiguity and component misuse
\end{itemize}

\subsection{4. Surface Selector Model}\label{record-042-surface-selector-model}

This makes stories and scenarios targetable in a stable way.

What it is:

\begin{itemize}
\tightlist
\item
  semantic targeting over surfaces and layout output
\item
  ids, classes, scoped regions, and text-based hooks where appropriate
\end{itemize}

Why it matters:

\begin{itemize}
\tightlist
\item
  docs and tests should not depend on brittle coordinates
\item
  selector semantics are the bridge between human-facing demos and machine-driven assertions
\end{itemize}

Primary assumption to test:

\begin{itemize}
\tightlist
\item
  a small selector model can cover most docs/demo/test needs without turning into a fake DOM
\end{itemize}

\subsection{5. Surface Scenario Test Harness}\label{record-042-surface-scenario-test-harness}

This turns stories into interactive truth surfaces.

What it is:

\begin{itemize}
\tightlist
\item
  a Playwright-like harness for Bijou apps and surface stories
\item
  scripted input, semantic selection, frame assertions, deterministic flows
\end{itemize}

Why it matters:

\begin{itemize}
\tightlist
\item
  it can make the same stories useful for docs demos and regression checks
\item
  it reduces the cost of verifying behavior that is currently trapped inside bespoke examples
\end{itemize}

Primary assumption to test:

\begin{itemize}
\tightlist
\item
  once stories and selectors exist, scenarios can become the shared layer for docs, bug repros, and regression testing
\end{itemize}

\subsection{6. Component Lab Mode}\label{record-042-component-lab-mode}

This is the interactive workbench inside the docs system, not a separate competing product.

What it is:

\begin{itemize}
\tightlist
\item
  a knob-driven or control-driven exploration surface for stories
\item
  state toggles, width toggles, profile toggles, maybe timing or pointer toggles
\end{itemize}

Why it matters:

\begin{itemize}
\tightlist
\item
  it helps authors and users poke at a component family without writing another one-off demo
\item
  but it should live inside the docs/story system, not alongside it as another disconnected example app
\end{itemize}

Primary assumption to test:

\begin{itemize}
\tightlist
\item
  a lab mode is most valuable when it is powered by stories, not by custom per-example controls
\end{itemize}

\section{The next pass}\label{record-042-the-next-pass}

The next pass should not try to build all six pieces at once. It should validate the smallest coherent loop:

\begin{enumerate}
\def\labelenumi{\arabic{enumi}.}
\tightlist
\item
  \textbf{Story Protocol v0}
\item
  \textbf{Docs TUI vertical slice}
\item
  \textbf{Interaction Profile Simulator}
\item
  \textbf{Surface Selector Model}
\item
  \textbf{Surface Scenario Test Harness}
\item
  \textbf{Component Lab Mode}
\item
  \textbf{Shell-owned settings drawer}
\end{enumerate}

\subsection{Why this order}\label{record-042-why-this-order}

\subsubsection{1. Story Protocol v0}\label{record-042-story-protocol-v0}

Build first because everything else gets cheaper and truer once stories exist.

The goal is not to design the perfect eternal schema. The goal is to prove that one story record can power:

\begin{itemize}
\tightlist
\item
  one docs page
\item
  one live demo
\item
  one source snippet
\item
  one profile comparison
\end{itemize}

Recommended first families:

\begin{itemize}
\tightlist
\item
  \texttt{a\allowbreak{}l\allowbreak{}e\allowbreak{}r\allowbreak{}t\allowbreak{}(\allowbreak{})} / \texttt{n\allowbreak{}o\allowbreak{}t\allowbreak{}e\allowbreak{}(\allowbreak{})} for simple lowering
\item
  \texttt{m\allowbreak{}o\allowbreak{}d\allowbreak{}a\allowbreak{}l\allowbreak{}(\allowbreak{})} / \texttt{d\allowbreak{}r\allowbreak{}a\allowbreak{}w\allowbreak{}e\allowbreak{}r\allowbreak{}(\allowbreak{})} for richer structured overlays
\item
  \texttt{n\allowbreak{}a\allowbreak{}v\allowbreak{}i\allowbreak{}g\allowbreak{}a\allowbreak{}b\allowbreak{}l\allowbreak{}e\allowbreak{}T\allowbreak{}a\allowbreak{}b\allowbreak{}l\allowbreak{}e\allowbreak{}S\allowbreak{}u\allowbreak{}r\allowbreak{}f\allowbreak{}a\allowbreak{}c\allowbreak{}e\allowbreak{}(\allowbreak{})} or \texttt{v\allowbreak{}i\allowbreak{}e\allowbreak{}w\allowbreak{}p\allowbreak{}o\allowbreak{}r\allowbreak{}t\allowbreak{}S\allowbreak{}u\allowbreak{}r\allowbreak{}f\allowbreak{}a\allowbreak{}c\allowbreak{}e\allowbreak{}(\allowbreak{})} for structured TUI behavior
\end{itemize}

\subsubsection{2. Docs TUI vertical slice}\label{record-042-docs-tui-vertical-slice}

Once the story shape exists, build a narrow docs app that renders a few canonical stories.

The first slice should include:

\begin{itemize}
\tightlist
\item
  left nav
\item
  content pane
\item
  demo pane
\item
  one or two width presets
\item
  one or two family pages
\end{itemize}

Do not try to build the full docs universe first.

\subsubsection{3. Interaction Profile Simulator}\label{record-042-interaction-profile-simulator-1}

Add the control that most directly demonstrates Bijou's difference.

The first simulator should support:

\begin{itemize}
\tightlist
\item
  \texttt{r\allowbreak{}i\allowbreak{}c\allowbreak{}h}
\item
  \texttt{s\allowbreak{}t\allowbreak{}a\allowbreak{}t\allowbreak{}i\allowbreak{}c}
\item
  \texttt{p\allowbreak{}i\allowbreak{}p\allowbreak{}e}
\item
  \texttt{a\allowbreak{}c\allowbreak{}c\allowbreak{}e\allowbreak{}s\allowbreak{}s\allowbreak{}i\allowbreak{}b\allowbreak{}l\allowbreak{}e}
\end{itemize}

Optional next toggles:

\begin{itemize}
\tightlist
\item
  reduced motion
\item
  narrow width
\item
  no color
\end{itemize}

\subsubsection{4. Surface Selector Model}\label{record-042-surface-selector-model-1}

Once stories and docs exist, define the semantic targeting needed to make them testable and inspectable.

Keep the first version minimal:

\begin{itemize}
\tightlist
\item
  \texttt{i\allowbreak{}d}
\item
  \texttt{c\allowbreak{}l\allowbreak{}a\allowbreak{}s\allowbreak{}s}
\item
  scoped region / pane identity
\item
  text contains
\end{itemize}

\subsubsection{5. Surface Scenario Test Harness}\label{record-042-surface-scenario-test-harness-1}

After stories and selectors exist, use them to prove the docs surface can also be exercised mechanically.

The first version should be able to:

\begin{itemize}
\tightlist
\item
  open a story
\item
  toggle profile or width
\item
  assert on semantic targets
\item
  capture a frame sequence if needed
\end{itemize}

\subsubsection{6. Component Lab Mode}\label{record-042-component-lab-mode-1}

Only after the above loop exists should the team add richer controls.

That keeps the lab from becoming another disconnected example swamp.

\section{Success criteria}\label{record-042-success-criteria}

DOGFOOD is successful when:

\begin{itemize}
\tightlist
\item
  a new contributor can learn a component family from one canonical surface
\item
  docs, demos, and regressions are visibly converging instead of drifting
\item
  at least one family can be documented, previewed, profile-switched, and scenario-tested from the same story record
\item
  the docs app feels like a real product surface, not a pile of examples behind a menu
\end{itemize}

\section{Non-goals for the first iteration}\label{record-042-non-goals-for-the-first-iteration}

\begin{itemize}
\tightlist
\item
  replacing every example immediately
\item
  building the full visual studio
\item
  perfecting every replay or selector capability before the docs TUI exists
\item
  inventing a giant story schema before proving a narrow vertical slice
\end{itemize}

\section{Recommended first slice}\label{record-042-recommended-first-slice}

If the team wants the most leverage with the least risk, start here:

\begin{enumerate}
\def\labelenumi{\arabic{enumi}.}
\tightlist
\item
  define \texttt{C\allowbreak{}o\allowbreak{}m\allowbreak{}p\allowbreak{}o\allowbreak{}n\allowbreak{}e\allowbreak{}n\allowbreak{}t\allowbreak{}S\allowbreak{}t\allowbreak{}o\allowbreak{}r\allowbreak{}y} v0
\item
  author three canonical stories:

  \begin{itemize}
  \tightlist
  \item
    \texttt{a\allowbreak{}l\allowbreak{}e\allowbreak{}r\allowbreak{}t\allowbreak{}(\allowbreak{})} / \texttt{n\allowbreak{}o\allowbreak{}t\allowbreak{}e\allowbreak{}(\allowbreak{})}
  \item
    \texttt{m\allowbreak{}o\allowbreak{}d\allowbreak{}a\allowbreak{}l\allowbreak{}(\allowbreak{})} / \texttt{d\allowbreak{}r\allowbreak{}a\allowbreak{}w\allowbreak{}e\allowbreak{}r\allowbreak{}(\allowbreak{})}
  \item
    \texttt{v\allowbreak{}i\allowbreak{}e\allowbreak{}w\allowbreak{}p\allowbreak{}o\allowbreak{}r\allowbreak{}t\allowbreak{}S\allowbreak{}u\allowbreak{}r\allowbreak{}f\allowbreak{}a\allowbreak{}c\allowbreak{}e\allowbreak{}(\allowbreak{})} or \texttt{n\allowbreak{}a\allowbreak{}v\allowbreak{}i\allowbreak{}g\allowbreak{}a\allowbreak{}b\allowbreak{}l\allowbreak{}e\allowbreak{}T\allowbreak{}a\allowbreak{}b\allowbreak{}l\allowbreak{}e\allowbreak{}S\allowbreak{}u\allowbreak{}r\allowbreak{}f\allowbreak{}a\allowbreak{}c\allowbreak{}e\allowbreak{}(\allowbreak{})}
  \end{itemize}
\item
  build a docs TUI shell that renders those stories
\item
  add profile switching to those pages
\end{enumerate}

That is the smallest slice that proves the philosophy instead of merely describing it.

